\documentclass[10pt,a4paper]{amsart}
\usepackage[margin=19mm]{geometry}
\usepackage{amsmath,amssymb,mathtools}
\usepackage[scr=rsfs,cal=euler]{mathalfa}
\usepackage{xfrac}
\usepackage[unicode,hidelinks,bookmarks=false]{hyperref}
\newcommand{\ie}{i.e.\ }
\newcommand{\cf}{cf.\ }
\newcommand{\resp}{respectively\ }
\newcommand{\Subsection}{\subsection}
\providecommand{\nobreakdash}{\nobreak\hbox{-}}
\setlength{\emergencystretch}{2em}
\multlinegap0pt
\usepackage[shortlabels]{enumitem}
\usepackage{mathtools}
\usepackage{amssymb}
\usepackage{xcolor}
\newtheorem{proposition}{Proposition}
\newcommand{\rhoCab}{{\rho_{\mathrm{cab}}}}
\title{Pr\"ufer $2$-group and Milnor's Conjecture on Fundamental Groups}
\author{Nan Wu, Zetian Yan}
\date{Source: arXiv:2608.15505v1 (CC BY 4.0)}
\begin{document}
\maketitle

\noindent\emph{Source and licence.} Pr\"ufer $2$-group and Milnor's Conjecture on Fundamental Groups. Version arXiv:2608.15505v1 (CC BY 4.0), distributed by arXiv under the Creative Commons Attribution 4.0 International licence, http://creativecommons.org/licenses/by/4.0/, as recorded in the arXiv licence metadata for this exact version and shown on the abstract page https://arxiv.org/abs/2608.15505v1. Full licence notice: \texttt{licenses/arxiv-2608.15505-CC-BY-4.0.txt}.\par\smallskip

\noindent\textit{Editorial scope.} Counted unit: the article's proposition prop:double-cable (source0.tex lines 800--848). The article's complete original proof of it is delivered with it (source0.tex lines 850--1127, one \texttt{proof} environment of 278 lines). No mathematics is written, repaired or re-derived: the statement and the proof are the article's own lines, in the article's own notation, and no retained line is substituted or re-flowed.  Retained uncounted as the source-local context the proof needs: lines 725--738.  Nothing was omitted from any retained span, so the package carries no omission record.  No retained line contains a citation, so the wrapper carries no bibliography and no bibliography entry.  No retained line contains a figure environment, an image-inclusion command or drawing code, so no image is delivered and the package declares no asset dependency.  Editorial formatting only, in addition: an AMS \texttt{amsart} preamble that restates the article's own declarations and macros used by the retained lines, namely the environments \texttt{proposition} on separate counters, together with the standard packages those lines need.  The retained environments print in this document as Proposition 1 in order of appearance, whereas the article prints the same items with the numbering of its own full document.  The complete native source of the article as distributed in the arXiv e-print for this exact version is bundled unchanged as \texttt{sources/207752/source0.tex}.
\begin{align}
 c_2(\theta)
   &:=\bigl(e^{2\mathrm i\theta},\varepsilon u(\theta)\bigr),\label{eq:double-cable}\\
 \mathcal F_{\rhoCab}(\theta,v)
   &:=\bigl(e^{2\mathrm i\theta},\varepsilon u(\theta)+v\bigr),\label{eq:explicit-tube}\\
 \nu_{\rhoCab}(c_2)
   &:=\mathcal F_{\rhoCab}
     \bigl((\mathbb R/2\pi\mathbb Z)\times
                     D^{m+1}_{\rhoCab}\bigr),\notag\\
 C_m
   &:=V\setminus\mathcal F_{\rhoCab}
     \bigl((\mathbb R/2\pi\mathbb Z)\times
              \operatorname{int}D^{m+1}_{\rhoCab}\bigr).\label{eq:explicit-exterior}
\end{align}

\begin{proposition}
\label{prop:double-cable}
The exterior $C_m$ defined above has the following properties.
\begin{enumerate}[label=\textup{(\roman*)}]
\item $\mathcal F_{\rhoCab}$ is a smooth embedding,
\[
 \partial C_m=\partial_oC_m\sqcup\partial_iC_m,
 \qquad
 \partial_oC_m\cong\mathbb S^1\times\mathbb S^m,
 \qquad
 \partial_iC_m\cong\mathbb S^1\times\mathbb S^m.
\]
\item If $j:C_m\hookrightarrow V$ is the inclusion, then
\[
 j_*:\pi_1(C_m,x_O)\xrightarrow{\ \cong\ }\pi_1(V,x_O),
 \qquad
 \pi_1(C_m,x_O)=\langle g\rangle\cong\mathbb Z,
 \qquad
 g=[\lambda_O].
\]
\item The path $\beta$ is contained in $C_m$, satisfies
\[
 \beta(0)=x_O,
 \qquad
 \beta(1)=x_I,
\]
and, with $\beta^{-1}(a)=\beta(1-a)$,
\[
 \bigl[\beta*\lambda_I*\beta^{-1}\bigr]=g^2
 \quad\text{in }\pi_1(C_m,x_O).
\]
For every path $\widehat\beta:[0,1]\to C_m$ from $x_O$ to $x_I$,
\[
 \bigl[\widehat\beta*\lambda_I*\widehat\beta^{-1}\bigr]=g^2.
\]
\item The map
\[
 \Phi:T_\phi\longrightarrow C_m,
 \qquad
 \Phi([x,\varphi])=(e^{\mathrm i\varphi},R_{\varphi/2}x),
\]
is a diffeomorphism, and
\[
 \phi(I_+)=I_-,
 \qquad
 \phi(I_-)=I_+.
\]
\end{enumerate}
\end{proposition}

\begin{proof}
\emph{Part \textup{(i)}.}
For $(\theta,v),(\theta',v')\in
(\mathbb R/2\pi\mathbb Z)\times D^{m+1}_{\rhoCab}$, the equality
$\mathcal F_{\rhoCab}(\theta,v)
 =\mathcal F_{\rhoCab}(\theta',v')$ gives
\[
 e^{2\mathrm i\theta}=e^{2\mathrm i\theta'},
 \qquad
 \theta'-\theta\in\{0,\pi\}\pmod{2\pi}.
\]
If $\theta'=\theta$, then $v'=v$.  If
$\theta'=\theta+\pi$, then
\[
 u(\theta')=-u(\theta),
 \qquad
 v'-v=2\varepsilon u(\theta),
\]
and hence
\[
 2\varepsilon=|v'-v|\le |v'|+|v|\le2\rhoCab<2\varepsilon,
\]
a contradiction.  Thus $\mathcal F_{\rhoCab}$ is injective.  For
$a\in\mathbb R$ and $w\in\mathbb R^{m+1}$, its differential is
\[
 d\mathcal F_{\rhoCab,(\theta,v)}(a,w)
 =\bigl(2a\mathrm i e^{2\mathrm i\theta},
         \varepsilon a u'(\theta)+w\bigr),
\]
and
\[
 d\mathcal F_{\rhoCab,(\theta,v)}(a,w)=0
 \quad\Longrightarrow\quad a=0\quad\Longrightarrow\quad w=0.
\]
Consequently
\[
 \operatorname{rank}d\mathcal F_{\rhoCab,(\theta,v)}=m+2.
\]
Since the domain of $\mathcal F_{\rhoCab}$ is compact and $V$ is Hausdorff,
$\mathcal F_{\rhoCab}$ is a smooth embedding.  The estimate
\[
 |\varepsilon u(\theta)+v|\le\varepsilon+\rhoCab<1
\]
gives
\[
 \nu_{\rhoCab}(c_2)\subset\mathbb S^1\times\operatorname{int}D^{m+1}.
\]
It follows from \eqref{eq:explicit-exterior} that
\[
 \partial C_m
 =\bigl(\mathbb S^1\times\partial D^{m+1}\bigr)
  \sqcup
  \mathcal F_{\rhoCab}
  \bigl((\mathbb R/2\pi\mathbb Z)\times
                         \partial D^{m+1}_{\rhoCab}\bigr)
 =\partial_oC_m\sqcup\partial_iC_m,
\]
and
\[
 \partial_oC_m\cong\mathbb S^1\times\mathbb S^m,
 \qquad
 \partial_iC_m\cong\mathbb S^1\times\mathbb S^m.
\]

\emph{Part \textup{(ii)}.}
Define
\[
\mathcal D_\tau:V\setminus\operatorname{im}c_2\longrightarrow
                  V\setminus\operatorname{im}c_2,
 \qquad 0\le\tau\le1,
\]
by
\[
\mathcal D_\tau(y)=
 \begin{cases}
 \displaystyle
 \mathcal F_{\rhoCab}\left(\theta,
  \left((1-\tau)|v|+\tau\rhoCab\right)\frac{v}{|v|}\right),
 &y=\mathcal F_{\rhoCab}(\theta,v),\quad0<|v|\le\rhoCab,\\[3mm]
 y,&y\in C_m.
 \end{cases}
\]
On $|v|=\rhoCab$ the two formulas agree, so $\mathcal D_\tau$ is
continuous.  Hence $\mathcal D_\tau$ is a strong deformation retraction
of $V\setminus\operatorname{im}c_2$ onto $C_m$; explicitly,
\begin{align*}
 \mathcal D_0&=\operatorname{id}_{V\setminus\operatorname{im}c_2},\\
 \mathcal D_\tau|_{C_m}&=\operatorname{id}_{C_m},\\
 \mathcal D_1\bigl(V\setminus\operatorname{im}c_2\bigr)&=C_m.
\end{align*}

Let $[\gamma]\in\pi_1(V,x_O)$.  Since $x_O\notin\operatorname{im}c_2$,
relative transversality gives
$\gamma'\simeq\gamma$ relative to $x_O$ with
$\gamma'\pitchfork\operatorname{im}c_2$.  Therefore
\[
 \dim (\gamma')^{-1}\bigl(\operatorname{im}c_2\bigr)
 =1+1-(m+2)=-m<0,
\]
so
\[
 (\gamma')^{-1}\bigl(\operatorname{im}c_2\bigr)=\varnothing,
 \qquad
 \mathcal D_1\circ\gamma':\mathbb S^1\longrightarrow C_m.
\]
Thus $j_*$ is surjective.

If $[\gamma]\in\ker j_*$, there is a map
\[
 H:(D^2,\partial D^2)\longrightarrow(V,C_m),
 \qquad H|_{\partial D^2}=\gamma.
\]
The boundary of $H$ is contained in $C_m$ and is therefore disjoint from
$\operatorname{im}c_2$.  Relative transversality gives
$H'\simeq H$ relative to $\partial D^2$ with
$H'\pitchfork\operatorname{im}c_2$.  Hence
\[
 \dim (H')^{-1}\bigl(\operatorname{im}c_2\bigr)
 =2+1-(m+2)=1-m<0,
\]
and therefore
\[
 (H')^{-1}\bigl(\operatorname{im}c_2\bigr)=\varnothing.
\]
The map $\mathcal D_1\circ H':D^2\to C_m$ satisfies
\[
 (\mathcal D_1\circ H')|_{\partial D^2}=\gamma,
\]
so $[\gamma]=1$ in $\pi_1(C_m,x_O)$.  Thus $j_*$ is injective.
Since
\[
 V\simeq\mathbb S^1,
 \qquad
 \pi_1(V,x_O)\cong\mathbb Z,
\]
we obtain
\[
 j_*:\pi_1(C_m,x_O)\xrightarrow{\ \cong\ }\pi_1(V,x_O),
 \qquad
 \pi_1(C_m,x_O)\cong\mathbb Z.
\]

\emph{Part \textup{(iii)}.}
The equation $e^{2\mathrm i\theta}=1$ has the two solutions
$\theta=0,\pi$ modulo $2\pi$.  Consequently,
\[
 \nu_{\rhoCab}(c_2)\cap(\{1\}\times D^{m+1})
 =\{1\}\times(D^{m+1}_+\sqcup D^{m+1}_-).
\]
Write
\[
 b(a)=a\varepsilon e_1+(1-a+a\rhoCab)e_{m+1},
 \qquad \beta(a)=(1,b(a)).
\]
Then
\[
 |b(a)|\le(1-a)+a(\varepsilon+\rhoCab)\le1,
\]
and
\begin{align*}
 |b(a)-\varepsilon e_1|^2
 &=\varepsilon^2(1-a)^2+(1-a+a\rhoCab)^2
 \ge\rhoCab^2,\\
 |b(a)+\varepsilon e_1|
 &\ge(1+a)\varepsilon>\rhoCab.
\end{align*}
Therefore
\[
 \beta([0,1])\subset C_m,
 \qquad
 \beta(0)=x_O,
 \qquad
 \beta(1)=x_I.
\]

Let
\[
 p_V:V\longrightarrow\mathbb S^1,
 \qquad p_V(z,x)=z,
 \qquad p=p_V|_{C_m}.
\]
Since $p_*=(p_V)_*\circ j_*$ and
$(p_V)_*:\pi_1(V,x_O)\to\pi_1(\mathbb S^1,1)$ is an isomorphism,
\[
 p_*:\pi_1(C_m,x_O)\xrightarrow{\ \cong\ }\mathbb Z.
\]
The identities
\begin{align*}
 (p\circ\lambda_O)(\varphi)&=e^{\mathrm i\varphi},\\
 p\circ\beta&\equiv1,\\
 (p\circ\lambda_I)(\theta)&=e^{2\mathrm i\theta}
\end{align*}
give
\begin{align*}
 p_*[\lambda_O]&=1,\\
 p_*\bigl[\beta*\lambda_I*\beta^{-1}\bigr]&=2.
\end{align*}
Consequently, for $g=[\lambda_O]$,
\[
 \bigl[\beta*\lambda_I*\beta^{-1}\bigr]=g^2.
\]
If $\widehat\beta$ is another path from $x_O$ to $x_I$, then
\[
 \bigl[\widehat\beta*\lambda_I*\widehat\beta^{-1}\bigr]
 =
 [\widehat\beta*\beta^{-1}]\,
 \bigl[\beta*\lambda_I*\beta^{-1}\bigr]\,
 [\widehat\beta*\beta^{-1}]^{-1}.
\]
The group $\pi_1(C_m,x_O)\cong\mathbb Z$ is abelian, so the right-hand
side equals $g^2$.

\emph{Part \textup{(iv)}.}
For $0\le\varphi\le2\pi$, the two solutions of
$e^{2\mathrm i\theta}=e^{\mathrm i\varphi}$ are
\[
 \theta=\frac\varphi2,
 \qquad
 \theta=\frac\varphi2+\pi
 \pmod{2\pi}.
\]
Since
\[
 \varepsilon u(\varphi/2)=R_{\varphi/2}(\varepsilon e_1),
 \qquad
 \varepsilon u(\varphi/2+\pi)=R_{\varphi/2}(-\varepsilon e_1),
\]
we have
\begin{align}
 \nu_{\rhoCab}(c_2)\cap
 \bigl(\{e^{\mathrm i\varphi}\}\times D^{m+1}\bigr)
 &=\{e^{\mathrm i\varphi}\}\times
 R_{\varphi/2}\bigl(D^{m+1}_+\sqcup D^{m+1}_-\bigr),
                                                   \label{eq:double-cable-tube-fiber}\\
 C_m\cap\bigl(\{e^{\mathrm i\varphi}\}\times D^{m+1}\bigr)
 &=\{e^{\mathrm i\varphi}\}\times R_{\varphi/2}P_m.
                                                   \label{eq:double-cable-exterior-fiber}
\end{align}
Furthermore,
\begin{align*}
 R_\pi(D^{m+1}_+)&=D^{m+1}_-,
 &R_\pi(D^{m+1}_-)&=D^{m+1}_+,\\
 \phi(I_+)&=I_-,
 &\phi(I_-)&=I_+.
\end{align*}
The map
\[
 \widetilde\Phi:P_m\times\mathbb R\longrightarrow C_m,
 \qquad
 \widetilde\Phi(x,\varphi)
   =(e^{\mathrm i\varphi},R_{\varphi/2}x),
\]
satisfies
\[
 \widetilde\Phi(x,\varphi+2\pi)
 =(e^{\mathrm i\varphi},R_{\varphi/2}R_\pi x)
 =\widetilde\Phi(\phi(x),\varphi).
\]
It therefore induces
\[
 \Phi:T_\phi\longrightarrow C_m,
 \qquad
 \Phi([x,\varphi])
   =(e^{\mathrm i\varphi},R_{\varphi/2}x).
\]
For fixed $\varphi$, the differential in the $P_m$-directions is the
isomorphism induced by $R_{\varphi/2}$, whereas the differential in the
$\varphi$-direction has nonzero component tangent to the first circle
factor.  Thus $\widetilde\Phi$, and hence $\Phi$, is a local
diffeomorphism.  Equation~\eqref{eq:double-cable-exterior-fiber} shows
that for $0\le\varphi<2\pi$ every point
$(e^{\mathrm i\varphi},y)\in C_m$ has the unique preimage
\[
 [R_{-\varphi/2}y,\varphi]\in T_\phi.
\]
Therefore $\Phi$ is bijective.  A bijective local diffeomorphism is a
diffeomorphism.
\end{proof}

\end{document}
